% ============================================================================
% Fichier  : partie2/P02_C06_PIU_processus_universel.tex
% Titre    : Processus Investigatif Universel
% Auteur   : MINKA MI NGUIDJOÏ Thierry Emmanuel
% Version  : 1.0 -- Juin 2026
% Statut   : RÉDIGÉ
% Source   : cours_IN_METHODOLOGIE.pdf (39 pages) -- fidèlement converti
%            et habillé selon le style TPIN. Contenu original intégralement
%            préservé. Enrichissements : encadrés CRO, boxPNC, ancrage
%            camerounais, liens avec les affaires MBONGO et ROUSSEAU.
% ============================================================================

\chapter{Processus Investigatif Universel}
\label{chap:PIU}

\epigraph{%
    « La qualité d'une investigation se mesure moins à la rapidité avec
    laquelle on parvient à des conclusions qu'à la robustesse avec laquelle
    ces conclusions résistent à la contestation contradictoire. »%
}{--- Principe fondamental du PIU}

\niveauChapitre{Fondamental à Intermédiaire}{%
    Lecture de la Partie~I recommandée. Ce chapitre constitue le c\oe ur
    méthodologique de ce cours~; toutes les parties suivantes s'y réfèrent.%
}

\begin{boxobjectifs}
\begin{itemize}
    \item Comprendre l'architecture modulaire du \termecle{Processus
          Investigatif Universel} (PIU)~: 5~phases, 17~étapes, 3~Quality Gates.
    \item Maîtriser la \termecle{matrice de criticité} pour calibrer le niveau
          de rigueur adapté à chaque investigation.
    \item Appliquer chaque étape selon trois niveaux de profondeur (minimal,
          standard, avancé) selon les contraintes du contexte.
    \item Conduire une investigation complète dans trois contextes distincts~:
          particulier, entreprise et judiciaire.
    \item Comprendre le rôle des Quality Gates comme mécanismes de contrôle
          qualité obligatoires.
    \item Relier chaque phase du PIU aux dimensions du Trilemme CRO.
\end{itemize}
\end{boxobjectifs}

\bigskip

% ============================================================================
\section{Introduction et Architecture Générale}
\label{sec:piu-intro}
% ============================================================================

\subsection{Objectifs et principes fondateurs}

L'investigation, qu'elle soit menée dans un cadre personnel, professionnel ou
judiciaire, nécessite une méthodologie rigoureuse garantissant à la fois
l'intégrité probatoire et la défendabilité des conclusions. Ce chapitre
présente un \termecle{Processus Investigatif Universel} (PIU), conçu pour
s'adapter à tout contexte investigatif par le biais d'une architecture
modulaire et scalable.

Le PIU vise quatre objectifs fondamentaux~:
\begin{enumerate}
    \item \textbf{Établissement factuel objectif}~: documenter les faits de
          manière reproductible et vérifiable~;
    \item \textbf{Intégrité probatoire}~: garantir que les éléments collectés
          conservent leur valeur probante~;
    \item \textbf{Traçabilité complète}~: permettre l'audit de chaque étape
          du processus investigatif~;
    \item \textbf{Défendabilité}~: produire des conclusions résistant à la
          contestation contradictoire.
\end{enumerate}

\subsection{Principe d'adaptabilité~: trois contextes}

Le PIU emploie une terminologie neutre permettant son application dans trois
contextes principaux~:
\begin{itemize}
    \item \textbf{Contexte particulier}~: investigation menée par un individu
          pour ses intérêts propres~;
    \item \textbf{Contexte organisationnel}~: investigation interne à une
          entreprise ou institution~;
    \item \textbf{Contexte judiciaire}~: expertise mandatée par une autorité
          juridictionnelle.
\end{itemize}

Chaque étape comporte trois niveaux de profondeur (minimal, standard, avancé)
s'ajustant selon la criticité du cas, conformément à la matrice de criticité
du Tableau~\ref{tab:piu-criticite}.

\subsection{Matrice de criticité}

\begin{tableaucours}{p{3.5cm}p{2.8cm}p{2.8cm}p{2.8cm}}{%
    \tableauHeader{Critère} &
    \tableauHeader{Minimal} &
    \tableauHeader{Standard} &
    \tableauHeader{Avancé}
}
Enjeu financier & $<$ 10k & 10--100k & $>$ 100k \\
Risque pénal & Nul & Contraventionnel & Correctionnel/Criminel \\
Complexité technique & Faible & Moyenne & Haute \\
Nombre parties prenantes & $<$ 5 & 5--20 & $>$ 20 \\
Durée estimée & $<$ 1~semaine & 1~sem.~-- 3~mois & $>$ 3~mois \\
\end{tableaucours}
\vspace{-0.8em}
\begin{center}{\small\itshape Tableau~\ref{tab:piu-criticite} ---
Matrice de criticité et niveau de rigueur requis}\end{center}
\label{tab:piu-criticite}

\subsection{Architecture globale~: 5~phases, 17~étapes, 3~Quality Gates}

Le PIU s'articule en cinq phases séquentielles comportant 17~étapes
opérationnelles et trois points de contrôle qualité (\termecle{Quality Gates}).

\begin{figure}[H]
\centering
\begin{tikzpicture}[
    phase/.style={
        draw=CouleurPrimaire, fill=CouleurDef, rounded corners=4pt,
        minimum width=9.5cm, minimum height=0.85cm,
        font=\small\bfseries\color{CouleurPrimaire}, align=center
    },
    qg/.style={
        draw=CouleurAccent, fill=CouleurAlerte, rounded corners=3pt,
        minimum width=3.2cm, minimum height=0.7cm,
        font=\footnotesize\bfseries\color{CouleurBordAlerte}, align=center
    },
    fl/.style={-{Stealth[length=5pt]}, thick, color=CouleurPrimaire}
]
\node[phase] (p1) at (0, 0)   {PHASE 1 --- Initialisation et Cadrage (E1--E3)};
\node[qg]    (q1) at (5.5, -0.85) {QG1 --- Validation\\Pré-Opérationnelle};
\node[phase] (p2) at (0, -1.7) {PHASE 2 --- Acquisition et Préservation (E4--E6)};
\node[qg]    (q2) at (5.5, -2.55) {QG2 --- Intégrité\\Probatoire};
\node[phase] (p3) at (0, -3.4) {PHASE 3 --- Investigation Analytique (E7--E10)};
\node[qg]    (q3) at (5.5, -4.25) {QG3 --- Robustesse\\Analytique};
\node[phase] (p4) at (0, -5.1) {PHASE 4 --- Formalisation et Transmission (E11--E13)};
\node[phase] (p5) at (0, -6.0) {PHASE 5 --- Suivi Post-Transmission (E14--E17)};
% Flèches
\draw[fl] (p1) -- (p2);
\draw[fl] (p2) -- (p3);
\draw[fl] (p3) -- (p4);
\draw[fl] (p4) -- (p5);
% Boucle itérative
\draw[fl, dashed, color=CouleurSecondaire]
    (p3.east) -- ++(0.4,0) |- node[right, font=\tiny, color=CouleurSecondaire,
    pos=0.25] {Itération\\si nécessaire} (p2.east);
\end{tikzpicture}
\caption{Architecture globale du Processus Investigatif Universel (PIU)}
\label{fig:piu-architecture}
\end{figure}

\subsection{Principe des Quality Gates}

Les \termecle{Quality Gates} (QG) constituent des points de contrôle
obligatoires validant la conformité avant progression vers la phase suivante.
Leur structure s'articule autour de critères graduels~:
\begin{itemize}
    \item \textbf{Critères obligatoires}~: applicables à tous les contextes~;
    \item \textbf{Critères standard}~: pour cas moyennement complexes~;
    \item \textbf{Critères avancés}~: pour investigations à haute criticité.
\end{itemize}

Un échec à un Quality Gate impose soit une remédiation immédiate, soit un retour
aux étapes antérieures concernées. Trois décisions possibles~:
\begin{description}
    \item[\textcolor{CouleurBordLizbiz}{\textbf{VERT}}] Tous critères obligatoires
          validés~: progression vers la phase suivante.
    \item[\textcolor{CouleurAccent}{\textbf{ORANGE}}] Critère non bloquant
          manquant~: progression conditionnelle avec plan de régularisation.
    \item[\textcolor{CouleurBordAlerte}{\textbf{ROUGE}}] Critère bloquant non
          validé~: suspension jusqu'à mise en conformité complète.
\end{description}

% ============================================================================
\section{Phase~1~: Initialisation et Cadrage}
\label{sec:piu-phase1}
% ============================================================================

La phase d'initialisation détermine la recevabilité de la demande investigative
et établit les conditions opérationnelles de sa conduite.

\subsection{Étape E1 --- Réception et Évaluation Préliminaire}

\textbf{Objectif~:} Comprendre la nature de la demande et évaluer sa recevabilité
avant tout engagement de ressources.

\subsubsection*{Actions par niveau}

\textbf{Niveau minimal~:}
\begin{itemize}
    \item Réception de la demande d'investigation (formelle ou informelle)~;
    \item Identification du commanditaire et vérification de sa légitimité~;
    \item Compréhension sommaire des faits allégués~;
    \item Évaluation première de la faisabilité technique et juridique.
\end{itemize}

\textbf{Niveau standard~:} En complément du niveau minimal~:
\begin{itemize}
    \item Demande de précisions écrites formelles au commanditaire~;
    \item Vérification des autorisations légales préalables nécessaires~;
    \item Estimation préliminaire des ressources humaines, techniques et
          temporelles requises~;
    \item Identification des contraintes réglementaires applicables
          (Loi~2024/017, Convention de Budapest selon contexte).
\end{itemize}

\textbf{Niveau avancé~:} En complément des niveaux précédents~:
\begin{itemize}
    \item Analyse de risque multi-dimensionnelle (légal, réputationnel,
          technique, financier)~;
    \item Consultation de spécialistes pour évaluation de faisabilité~;
    \item Étude de précédents similaires (jurisprudence, cas antérieurs)~;
    \item Rédaction d'une note d'opportunité structurée.
\end{itemize}

\textbf{Livrable~:} \textit{Fiche d'Évaluation Initiale} comprenant~:
description factuelle de la demande~; enjeux identifiés (financiers, juridiques,
humains, réputationnels)~; obstacles prévisibles et risques~;
recommandation préliminaire motivée (accepter / refuser / clarifier).

\subsection{Étape E2 --- Décision d'Engagement}

\textbf{Objectif~:} Valider formellement le lancement de l'investigation selon
les critères de gouvernance appropriés au contexte.

\subsubsection*{Actions par niveau}

\textbf{Niveau minimal~:}
\begin{itemize}
    \item Décision personnelle d'engagement documentée~;
    \item Vérification d'absence de conflit d'intérêt personnel~;
    \item Validation de la disponibilité temporelle et des ressources.
\end{itemize}

\textbf{Niveau standard~:}
\begin{itemize}
    \item Validation hiérarchique ou contractuelle formalisée~;
    \item Signature d'un mandat d'investigation ou lettre de mission~;
    \item Définition précise du périmètre d'investigation et des limitations~;
    \item Établissement d'un budget prévisionnel et allocation de ressources~;
    \item Détermination des délais et jalons intermédiaires.
\end{itemize}

\textbf{Niveau avancé~:}
\begin{itemize}
    \item Validation collégiale impliquant expertises multiples (technique,
          juridique, éthique)~;
    \item Obtention des autorisations formelles requises (mandat judiciaire,
          réquisition~; art.~52 \loicam{2010/012} en contexte camerounais)~;
    \item Définition du niveau de classification sécuritaire~;
    \item Constitution d'un comité de pilotage si investigation complexe.
\end{itemize}

\textbf{Livrable~:} \textit{Mandat d'Investigation} comprenant~: périmètre précis
et exclusions explicites~; délais avec jalons intermédiaires~; budget et
ressources allouées~; niveau d'autorité conféré~; obligations de confidentialité~;
conditions de transmission finale.

\subsection{Étape E3 --- Planification Opérationnelle}

\textbf{Objectif~:} Structurer l'investigation en mobilisant les ressources
adéquates et en anticipant les difficultés opérationnelles.

\subsubsection*{Actions par niveau}

\textbf{Niveau minimal~:}
\begin{itemize}
    \item Auto-évaluation des compétences requises versus disponibles~;
    \item Listage des outils et équipements nécessaires~;
    \item Création d'un calendrier sommaire avec échéances principales.
\end{itemize}

\textbf{Niveau standard~:}
\begin{itemize}
    \item Constitution d'une équipe avec attribution claire des rôles~;
    \item Rédaction d'un plan d'investigation structuré (axes, méthodologies,
          jalons, hypothèses préliminaires)~;
    \item Configuration et test des outils techniques~;
    \item Mise en place d'un système de documentation et traçabilité.
\end{itemize}

\textbf{Niveau avancé~:}
\begin{itemize}
    \item Planification détaillée type gestion de projet (WBS, Gantt)~;
    \item Analyse de risque opérationnelle exhaustive avec plan de contingence~;
    \item Test et calibration des outils critiques~;
    \item Création d'un environnement de travail sécurisé~;
    \item Établissement de protocoles d'escalade pour situations critiques.
\end{itemize}

\textbf{Livrable~:} \textit{Plan d'Investigation v1.0} comprenant~: objectifs
et critères de succès mesurables~; méthodologie générale~; ressources mobilisées~;
échéancier prévisionnel~; procédures de gestion des imprévus.

\subsection{Quality Gate~1 --- Validation Pré-Opérationnelle}

\textbf{Objectif~:} Garantir que l'ensemble des conditions préalables au démarrage
effectif sont satisfaites.

\textbf{Critères obligatoires (tous contextes)~:}
\begin{itemize}
    \item Mandat d'investigation formalisé et validé par l'autorité compétente~;
    \item Périmètre clairement défini, compris et accepté par toutes parties~;
    \item Ressources minimales identifiées et effectivement disponibles~;
    \item Plan d'investigation v1.0 rédigé et approuvé~;
    \item Outils nécessaires opérationnels ou accessibles dans les délais.
\end{itemize}

\textbf{Critères standard~:}
\begin{itemize}
    \item Validation légale et/ou déontologique obtenue auprès des instances~;
    \item Protocoles de documentation établis et compris de tous~;
    \item Système de traçabilité configuré et testé.
\end{itemize}

\textbf{Critères avancés~:}
\begin{itemize}
    \item Plan de contingence documenté couvrant les risques identifiés~;
    \item Tests opérationnels des outils techniques réalisés~;
    \item Assurances et couvertures légales vérifiées si applicable.
\end{itemize}

\begin{boxPNC}
\textbf{Application OPJ --- Contexte camerounais, QG1~:}

En contexte judiciaire camerounais, le passage au VERT du QG1 exige~:
\begin{enumerate}
    \item Réquisition judiciaire ou mandat du procureur (art.~52 \loicam{2010/012})~;
    \item Disponibilité d'un kit forensique de base (write blocker,
          FTK Imager, clé USB bootable)~;
    \item Information du suspect de ses droits (art.~5 CPP Cameroun)
          si l'investigation implique une perquisition.
\end{enumerate}
\end{boxPNC}

% ============================================================================
\section{Phase~2~: Acquisition et Préservation}
\label{sec:piu-phase2}
% ============================================================================

La phase d'acquisition vise à constituer la base probatoire de l'investigation
en garantissant l'intégrité et la recevabilité juridique des éléments collectés.

\subsection{Étape E4 --- Collecte Systématique des Éléments Probatoires}

\textbf{Objectif~:} Rassembler exhaustivement les éléments factuels pertinents
en préservant leur intégrité et leur valeur probatoire.

\subsubsection*{Actions par niveau}

\textbf{Niveau minimal~:}
\begin{itemize}
    \item Identification systématique des sources d'information disponibles~;
    \item Collecte documentaire (contrats, correspondances, factures, notes)~;
    \item Réalisation de photographies ou copies des éléments physiques~;
    \item Rédaction de notes de constatation horodatées.
\end{itemize}

\textbf{Niveau standard~:}
\begin{itemize}
    \item Application d'un protocole de collecte standardisé et documenté~;
    \item Documentation de chaque action~: date/heure, méthode, outils,
          identité du collecteur, état initial, localisation~;
    \item Attribution d'une numérotation séquentielle unique à chaque élément~;
    \item Photographie de contexte (scène, emplacement) avant prélèvement.
\end{itemize}

\textbf{Niveau avancé~:}
\begin{itemize}
    \item Application stricte des protocoles \norme{ISO~27037} et
          principes ACPO~;
    \item Création d'images forensiques bit-à-bit avec write blocker~;
    \item Calcul d'empreintes cryptographiques (SHA-256 et SHA-512) pour
          tous éléments numériques~;
    \item Réalisation de prélèvements en présence de témoins assermentés
          si contexte judiciaire.
\end{itemize}

\textbf{Livrable~:} \textit{Inventaire des Éléments Probatoires} comprenant~:
liste numérotée~; description non ambiguë~; provenance et circonstances~;
état initial~; empreintes cryptographiques~; chaîne de traçabilité initiale.

\subsection{Étape E5 --- Sécurisation et Préservation}

\textbf{Objectif~:} Garantir l'intégrité durable des éléments collectés et
prévenir toute altération, perte ou accès non autorisé.

\subsubsection*{Actions par niveau}

\textbf{Niveau minimal~:}
\begin{itemize}
    \item Création immédiate de copies de sauvegarde sur support physique distinct~;
    \item Stockage sécurisé dans lieu physique protégé~;
    \item Règle stricte~: interdiction de travailler directement sur originaux~;
    \item Journal des accès (qui, quand, pourquoi).
\end{itemize}

\textbf{Niveau standard~:} Principe de sauvegarde 3-2-1~:
\begin{itemize}
    \item 3~copies de chaque élément~;
    \item 2~types de supports différents (disque dur, cloud, support optique)~;
    \item 1~copie en localisation géographique distincte~;
    \item Chiffrement systématique des données sensibles (AES-256)~;
    \item Vérification régulière (hebdomadaire) de l'intégrité par hash.
\end{itemize}

\textbf{Niveau avancé~:}
\begin{itemize}
    \item Utilisation de blockchain privée ou horodatage qualifié~;
    \item Stockage en coffre-fort numérique certifié~;
    \item Tests d'intégrité automatisés quotidiens~;
    \item Plan de continuité opérationnelle documenté.
\end{itemize}

\textbf{Livrable~:} \textit{Registre de Préservation} comprenant~:
certificats d'intégrité (hash avec algorithme)~; schéma de localisation des
copies~; journal d'accès horodaté~; résultats des tests d'intégrité.

\subsection{Étape E6 --- Identification et Cartographie des Parties Prenantes}

\textbf{Objectif~:} Recenser exhaustivement toutes les personnes et entités ayant
un lien direct ou indirect avec les faits investigués.

\subsubsection*{Actions par niveau}

\textbf{Niveau minimal~:}
\begin{itemize}
    \item Listage des parties directement mentionnées~;
    \item Distinction sommaire entre acteurs principaux et secondaires~;
    \item Recueil des coordonnées disponibles~;
    \item Notation des relations apparentes entre parties.
\end{itemize}

\textbf{Niveau standard~:}
\begin{itemize}
    \item Cartographie relationnelle détaillée (hiérarchiques, contractuels,
          familiaux)~;
    \item Classification par catégories~: parties directes~; témoins directs~;
          témoins indirects~; experts techniques~; autorités de tutelle~;
    \item Priorisation raisonnée pour la conduite des entretiens~;
    \item OSINT systématique dans limites légales (cf. Chapitre~\ref{chap:etat-art}).
\end{itemize}

\textbf{Niveau avancé~:}
\begin{itemize}
    \item Analyse de réseau social formelle (SNA) avec logiciels dédiés~;
    \item Profilage préliminaire adapté au contexte~;
    \item Évaluation du risque de collusion ou de coordination entre parties.
\end{itemize}

\textbf{Livrable~:} \textit{Matrice des Parties Prenantes} avec pour chaque
personne/entité~: identité et coordonnées~; rôle et catégorie~; niveau
d'implication estimé (1--5)~; crédibilité préliminaire~; priorité d'entretien~;
schéma relationnel visuel.

\subsection{Quality Gate~2 --- Intégrité Probatoire}

\textbf{Objectif~:} Vérifier que la base probatoire constituée est complète,
intègre, traçable et juridiquement exploitable.

\textbf{Critères obligatoires~:}
\begin{itemize}
    \item Inventaire probatoire exhaustif, numéroté de manière unique~;
    \item Chaîne de traçabilité documentée sans interruption pour chaque élément~;
    \item Sauvegardes multiples vérifiées et opérationnelles~;
    \item Impossibilité technique de modification non tracée~;
    \item Journal d'accès à jour, sans anomalie détectée.
\end{itemize}

\textbf{Critères standard~:}
\begin{itemize}
    \item Hash cryptographiques calculés et archivés dans localisation distincte~;
    \item Tests d'intégrité périodiques réalisés (100\% des hash correspondants)~;
    \item Conformité légale de toutes les méthodes de collecte validée.
\end{itemize}

\textbf{Rationale critique~:} Un élément probatoire non conforme à QG2 risque
l'irrecevabilité juridique (vice de procédure), la contestabilité (doute sur
authenticité) ou l'inexploitabilité (altération). La détection précoce évite
l'invalidation tardive de l'ensemble du travail investigatif.

\begin{boxCRO}
\textbf{Analyse CRO --- Phase~2~: Acquisition et Préservation}

$\mathcal{C}$~: La collecte exhaustive (E4 niveau avancé) peut inclure des
données sans rapport avec l'affaire. La proportionnalité doit guider le
périmètre de collecte.

$\mathcal{R}$~: Le principe 3-2-1 (E5) et le hachage SHA-256+SHA-512 maximisent
la fiabilité. L'impossibilité technique de modification non tracée est le
critère QG2 le plus important.

$\mathcal{O}$~: La chaîne de traçabilité continue et les hash documentés sont
les preuves d'intégrité opposables en justice.

\cro{$\downarrow$ proportionnalité}{$\uparrow\uparrow$ maximisée}{$\uparrow$ par la custody}
\end{boxCRO}

% ============================================================================
\section{Phase~3~: Investigation Analytique}
\label{sec:piu-phase3}
% ============================================================================

La phase d'investigation constitue le c\oe ur du processus, durant laquelle
hypothèses explicatives sont générées, testées et confrontées aux éléments
probatoires. Cette phase comporte une boucle itérative formalisée.

\subsection{Étape E7 --- Interrogatoires et Collecte Testimoniale}

\textbf{Objectif~:} Recueillir les témoignages, déclarations et expertises orales
nécessaires à la compréhension approfondie des faits.

\subsubsection*{Actions par niveau}

\textbf{Niveau minimal~:}
\begin{itemize}
    \item Conduite d'entretiens informels non structurés~;
    \item Prise de notes manuscrites contemporaines~;
    \item Recueil des déclarations spontanées~;
    \item Conservation soigneuse des témoignages écrits volontairement fournis.
\end{itemize}

\textbf{Niveau standard~:}
\begin{itemize}
    \item Préparation de questionnaires structurés adaptés à chaque catégorie~;
    \item Conduite d'entretiens semi-directifs (narration libre puis précisions)~;
    \item Enregistrement audio avec consentement formel écrit~;
    \item Rédaction de comptes rendus fidèles soumis à validation~;
    \item Analyse de cohérence interne (chronologie, contradictions)~;
    \item Documentation des refus de répondre.
\end{itemize}

\textbf{Niveau avancé~:}
\begin{itemize}
    \item Application de protocoles scientifiques~: méthode PEACE
          (\textit{Preparation, Engage, Account, Closure, Evaluation})~;
          entretien cognitif pour témoins de bonne foi~;
    \item Enregistrement multi-modal (audio et vidéo) si approprié~;
    \item Analyse linguistique approfondie des déclarations~;
    \item Procédures d'audition formelles respectant strictement les droits
          processuels (contexte judiciaire).
\end{itemize}

\textbf{Livrable~:} \textit{Corpus Testimonial} comprenant~: enregistrements
archivés~; transcriptions intégrales horodatées~; comptes rendus synthétiques~;
matrice d'évaluation de cohérence et crédibilité~; points de convergence et
divergences~; éléments corroborant ou réfutant les hypothèses.

\subsection{Étape E8 --- Génération d'Hypothèses Multiples}

\textbf{Objectif~:} Construire de manière systématique plusieurs scénarios
explicatifs distincts des faits, sans présupposé hiérarchique initial,
afin d'éviter le biais de confirmation.

\subsubsection*{Actions par niveau}

\textbf{Niveau minimal~:}
\begin{itemize}
    \item Formulation de 2 à 3 hypothèses principales spontanées~;
    \item Description narrative sommaire de chaque scénario~;
    \item Identification intuitive de l'hypothèse la plus probable.
\end{itemize}

\textbf{Niveau standard~:}
\begin{itemize}
    \item Génération systématique d'au moins 5~hypothèses distinctes incluant~:
          hypothèse principale, alternatives crédibles, hypothèse nulle
          (absence de fait délictuel), hypothèses improbables~;
    \item Documentation du raisonnement sous-jacent à chaque hypothèse~;
    \item Construction d'arbres de décision ou diagrammes causaux~;
    \item Formulation de prédictions testables dérivées
          (\textit{si H vraie, alors on devrait observer\ldots}).
\end{itemize}

\textbf{Niveau avancé~:}
\begin{itemize}
    \item Application formalisée de la méthode ACH
          (\textit{Analysis of Competing Hypotheses}) \cite{heuer1999}~;
    \item Modélisation bayésienne des probabilités a priori~;
    \item Exercice de \textit{red teaming}~: simulation d'acteur adversaire~;
    \item Analyse de précédents historiques similaires (benchmark).
\end{itemize}

\textbf{Livrable~:} \textit{Arbre d'Hypothèses v1.0} comprenant pour chaque
hypothèse~: identifiant unique et intitulé non ambigu~; description narrative~;
raisonnement explicite~; prédictions testables~; preuves pro/contra~;
conditions de réfutation~; probabilité subjective initiale.

\subsection{Étape E9 --- Confrontation Hypothèses/Preuves}

\textbf{Objectif~:} Évaluer systématiquement la compatibilité de chaque hypothèse
avec l'ensemble des éléments probatoires et témoignages collectés.

\subsubsection*{Actions par niveau}

\textbf{Niveau standard~:}
\begin{itemize}
    \item Construction d'une matrice de compatibilité hypothèses $\times$ preuves~;
    \item Pour chaque intersection, évaluation qualitative~:
          \textbf{+} Cohérente~; \textbf{-} Incohérente~; \textbf{o} Neutre~;
    \item Calcul de scores de cohérence par hypothèse~;
    \item Identification des preuves discriminantes~;
    \item Test empirique des prédictions dérivées (E8)~;
    \item Hiérarchisation des hypothèses par score de cohérence.
\end{itemize}

\textbf{Niveau avancé~:}
\begin{itemize}
    \item Analyse bayésienne formelle~;
    \item Simulation Monte-Carlo pour quantification des incertitudes~;
    \item Organisation de sessions de \textit{peer review}~;
    \item Processus de \textit{red team} cherchant activement à réfuter.
\end{itemize}

La figure suivante illustre un exemple de matrice de compatibilité~:

\begin{figure}[H]
\centering
\begin{tikzpicture}
\matrix (m) [matrix of nodes,
    nodes in empty cells,
    column sep=0.6cm, row sep=0.3cm,
    nodes={minimum width=1.2cm, minimum height=0.7cm,
           draw=CouleurBordInfo, font=\small, align=center}]
{
    {} & {H1} & {H2} & {H3} & {H4} \\
    {P1} & {\textcolor{red}{--}} & {\textcolor{CouleurBordLizbiz}{+}} &
            {\textcolor{CouleurBordLizbiz}{+}} & {o} \\
    {P2} & {o} & {\textcolor{CouleurBordLizbiz}{+}} &
            {\textcolor{red}{--}} & {\textcolor{CouleurBordLizbiz}{+}} \\
    {P3} & {\textcolor{CouleurBordLizbiz}{+}} & {o} &
            {\textcolor{CouleurBordLizbiz}{+}} & {o} \\
    {P4} & {\textcolor{red}{--}} & {\textcolor{red}{--}} & {o} &
            {\textcolor{CouleurBordLizbiz}{+}} \\
    {Score} & {1/4} & {2/4} & {2/4} & {3/4} \\
};
% Colorier la ligne d'en-tête
\foreach \col in {2,3,4,5} {
    \fill[CouleurPrimaire!20] (m-1-\col.north west) rectangle (m-1-\col.south east);
}
% Colorier la colonne d'en-tête
\foreach \row in {2,3,4,5,6} {
    \fill[CouleurPrimaire!20] (m-\row-1.north west) rectangle (m-\row-1.south east);
}
\end{tikzpicture}
\caption{Exemple de matrice de compatibilité hypothèses/preuves~:
\textbf{+} Cohérent~; \textbf{--} Incohérent~; \textbf{o} Neutre}
\label{fig:matrice-ach}
\end{figure}

\textbf{Livrable~:} \textit{Matrice d'Analyse Hypothèses/Preuves} comprenant~:
tableau croisé~; scores de cohérence~; liste des hypothèses réfutées~;
classement hiérarchique~; lacunes probatoires~; recommandations de collecte.

\subsection{Étape E9bis --- Falsification Active}

\textbf{Objectif~:} Appliquer le principe épistémologique de Popper~: tenter
activement de réfuter chaque hypothèse plutôt que de chercher uniquement des
confirmations.

\textbf{Justification méthodologique~:} Le biais de confirmation constitue l'un
des obstacles majeurs à l'objectivité investigative. La falsification active,
inspirée de la philosophie des sciences de Karl Popper \cite{popper1973},
impose la recherche délibérée d'observations qui, si découvertes, invalideraient
l'hypothèse testée. Une hypothèse ayant résisté à de multiples tentatives
sérieuses de falsification acquiert une robustesse épistémologique supérieure.

\subsubsection*{Actions par niveau}

\textbf{Niveau minimal~:}
\begin{itemize}
    \item Pour chaque hypothèse post-E9~: formulation de la question~: « Quelle
          observation prouverait de manière irréfutable que cette hypothèse est
          fausse~? »~;
    \item Vérification rapide dans le corpus probatoire~;
    \item Documentation sommaire.
\end{itemize}

\textbf{Niveau standard~:}
\begin{itemize}
    \item Pour chaque hypothèse viable, processus structuré en 4~étapes~:
          (1)~identification des observations critiques~;
          (2)~recherche active dans les preuves existantes~;
          (3)~planification de collecte additionnelle si non trouvées~;
          (4)~documentation des tentatives (succès ou échec)~;
    \item Calcul d'un indice de robustesse~: nombre de tentatives résistées.
\end{itemize}

\textbf{Niveau avancé~:}
\begin{itemize}
    \item Constitution d'une équipe \textit{Red Team} dédiée à la réfutation~;
    \item Analyse contrefactuelle~: « Si l'hypothèse était fausse,
          qu'observerait-on de différent~? »~;
    \item Formalisation logique des tests de falsification.
\end{itemize}

\textbf{Livrable~:} \textit{Rapport de Falsification} comprenant pour chaque
hypothèse testée~: conditions de falsification~; résultats~; score de robustesse~;
liste des hypothèses définitivement éliminées.

\subsection{Étape E10 --- Validation et Vérification}

\textbf{Objectif~:} Soumettre les conclusions préliminaires émergentes à un
processus rigoureux de validation multi-méthodes avant formalisation définitive.

\subsubsection*{Actions par niveau}

\textbf{Niveau minimal~:}
\begin{itemize}
    \item Relecture personnelle critique du raisonnement~;
    \item Vérification arithmétique des analyses quantitatives~;
    \item Auto-évaluation des biais cognitifs potentiels.
\end{itemize}

\textbf{Niveau standard~:}
\begin{itemize}
    \item Validation par pairs~: présentation structurée à collègues~;
    \item Test de cohérence globale~: narration complète et chronologique~;
    \item Identification et documentation des limitations méthodologiques~;
    \item Évaluation du niveau de certitude pour chaque conclusion.
\end{itemize}

\textbf{Niveau avancé~:}
\begin{itemize}
    \item Processus \textit{Red Team / Blue Team} formalisé~;
    \item Réexécution à l'aveugle par expert tiers~;
    \item Calcul d'intervalles de confiance statistiques~;
    \item Simulation d'argumentation adverse.
\end{itemize}

\textbf{Livrable~:} \textit{Rapport de Validation} comprenant~: méthodes de
validation employées~; objections identifiées et réponses~; limitations
reconnues~; niveau de consensus~; certification de validation (signatures).

\subsection{Quality Gate~3 --- Robustesse Analytique}

\textbf{Objectif~:} S'assurer que les conclusions sont suffisamment robustes,
documentées et défendables pour justifier leur formalisation.

\textbf{Critères obligatoires~:}
\begin{itemize}
    \item Génération documentée d'au moins 5~hypothèses distinctes (E8)~;
    \item Matrice de compatibilité hypothèses/preuves complétée (E9)~;
    \item Tentatives de falsification documentées pour hypothèses principales~;
    \item Application d'au moins une méthode de validation formelle (E10)~;
    \item Limitations de l'investigation clairement identifiées~;
    \item Traçabilité complète du raisonnement.
\end{itemize}

\textbf{Décision et gestion de l'itération~:}
\begin{description}
    \item[\textcolor{CouleurBordLizbiz}{\textbf{VERT}}] Progression vers Phase~4.
    \item[\textcolor{CouleurAccent}{\textbf{ORANGE}}] Expertise complémentaire
          ou re-analyse ciblée (5--10~jours).
    \item[\textcolor{CouleurBordAlerte}{\textbf{ROUGE}}] Retour selon nature~:
          lacunes testimoniales $\to$ E7~;
          incohérences analytiques $\to$ E8--E9~;
          preuves matérielles manquantes $\to$ E4--E6.
\end{description}

\begin{boiteRemarque}{Point critique sur l'itération}
Le retour en Phase~2 ou la réitération de la Phase~3 n'est pas un échec
méthodologique mais une \textbf{caractéristique normale} du processus
investigatif. Les investigations de qualité nécessitent fréquemment 2 à 3
passages dans la boucle analytique. Un taux de passage direct au QG3 de
60--70\% est satisfaisant (versus 85--90\% pour QG1 et QG2).
\end{boiteRemarque}

% ============================================================================
\section{Phase~4~: Formalisation et Transmission}
\label{sec:piu-phase4}
% ============================================================================

La phase de formalisation transforme le travail investigatif en livrables
structurés, adaptés au destinataire et juridiquement recevables.

\subsection{Étape E11 --- Documentation Exhaustive et Consolidation}

\textbf{Objectif~:} Rassembler, structurer et archiver l'intégralité des
éléments du dossier dans un format pérenne et auditable.

\textbf{Structure documentaire recommandée} --- le dossier consolidé
s'organise en cinq volumes~:
\begin{description}
    \item[\textbf{Volume~1 --- Synthèse}] Sommaire exécutif, chronologie
          factuelle, conclusions principales (10--20~pages).
    \item[\textbf{Volume~2 --- Analyses}] Méthodologie détaillée, hypothèses
          étudiées, raisonnement analytique (50--200~pages).
    \item[\textbf{Volume~3 --- Preuves}] Éléments probatoires numérotés avec
          certificats d'intégrité et chaînes de traçabilité.
    \item[\textbf{Volume~4 --- Témoignages}] Comptes rendus d'entretiens,
          déclarations écrites, expertises externes.
    \item[\textbf{Volume~5 --- Annexes}] Journal de bord, références,
          documents de travail.
\end{description}

\textbf{Livrable~:} \textit{Dossier d'Investigation Complet} structuré,
numéroté, indexé et archivé de manière pérenne.

\subsection{Étape E12 --- Rédaction du Livrable Final}

\textbf{Objectif~:} Produire un document de synthèse adapté au destinataire,
présentant de manière claire, rigoureuse et défendable les conclusions.

\subsubsection*{Structure recommandée du rapport}

\textbf{Niveau~1 --- Synthèse Exécutive (2--4~pages)~:}
\begin{enumerate}
    \item Contexte et objet de la mission~;
    \item Question(s) posée(s)~;
    \item Conclusions principales numérotées~;
    \item Niveau de confiance par conclusion (qualitatif ou quantitatif)~;
    \item Recommandations actionables immédiates~;
    \item Limitations assumées.
\end{enumerate}

\textbf{Niveau~2 --- Rapport Technique (20--100~pages)~:}
\begin{enumerate}
    \item Introduction et contexte approfondi~;
    \item Méthodologie (référencement normes/standards)~;
    \item Faits matériellement établis (chronologie détaillée)~;
    \item Hypothèses étudiées avec arbre décisionnel~;
    \item Analyse détaillée des preuves et témoignages~;
    \item Raisonnement étape par étape~;
    \item Incertitudes et limitations~;
    \item Recommandations détaillées~;
    \item Bibliographie.
\end{enumerate}

\textbf{Niveau~3 --- Dossier Probatoire (annexes)~:} Ensemble des pièces
justificatives référencées dans le rapport technique.

\subsubsection*{Principes rédactionnels fondamentaux}
\begin{description}
    \item[\textbf{Clarté}] Phrases courtes (15--20~mots maximum)~;
          vocabulaire précis~; structure logique avec transitions explicites.
    \item[\textbf{Objectivité}] Ton neutre (éviter « je pense »,
          « il est évident »)~; différenciation rigoureuse entre faits établis,
          interprétations et opinions~; limitations assumées.
    \item[\textbf{Défendabilité}] Chaque affirmation appuyée par référence
          probatoire~; raisonnement sans sauts inférentiels~;
          traçabilité complète du raisonnement.
    \item[\textbf{Adaptabilité}] Niveau de technicité ajusté selon
          destinataire~; synthèse exécutive pour décideurs non techniques.
\end{description}

\subsection{Étape E13 --- Transmission Sécurisée du Livrable}

\textbf{Objectif~:} Remettre le livrable de manière sécurisée, traçable,
conforme aux exigences de confidentialité.

\subsubsection*{Actions par niveau}

\textbf{Niveau minimal~:}
\begin{itemize}
    \item Remise en main propre contre signature ou email sécurisé chiffré~;
    \item Obtention d'un accusé de réception daté~;
    \item Conservation de la preuve de transmission.
\end{itemize}

\textbf{Niveau standard~:}
\begin{itemize}
    \item Transmission sécurisée~: chiffrement bout-en-bout (GPG, S/MIME)~;
    \item Bordereau de transmission détaillé et signé~;
    \item Liste exhaustive des documents transmis.
\end{itemize}

\textbf{Niveau avancé~:}
\begin{itemize}
    \item Remise en présence de témoins (huissier si contexte critique)~;
    \item Signature électronique qualifiée~;
    \item Horodatage blockchain pour preuve d'antériorité immuable~;
    \item Watermarking individualisé par destinataire.
\end{itemize}

\textbf{Livrable~:} \textit{Preuve de Transmission} comprenant bordereau signé,
liste exhaustive des documents, modalités de transmission, accusé contresigné.

% ============================================================================
\section{Phase~5~: Suivi Post-Transmission}
\label{sec:piu-phase5}
% ============================================================================

Cette phase, optionnelle pour investigations simples sans contentieux, devient
essentielle dans les contextes judiciaires ou lorsque des contestations sont
anticipées.

\subsection{Étape E14 --- Préparation à la Défense Contradictoire}

\textbf{Objectif~:} Anticiper les contestations possibles et préparer une défense
rigoureuse et documentée des conclusions.

\textbf{Actions~:} Constitution d'un dossier de défense comprenant~: analyse
des vulnérabilités du rapport~; arguments adverses anticipés avec contre-arguments~;
synthèses aide-mémoire par thème~; simulations d'interrogatoires contradictoires~;
références scientifiques ou jurisprudentielles de soutien.

\subsection{Étape E15 --- Maintenance Probatoire Continue}

\textbf{Objectif~:} Préserver intégrité et accessibilité du dossier durant toute
la procédure contentieuse ou administrative.

\textbf{Actions~:} Surveillance active (tests d'intégrité mensuels)~;
gestion rigoureuse des accès avec journal exhaustif~;
mise à jour de la chaîne de traçabilité~;
réponse aux demandes de clarification~;
correction d'erreurs matérielles avec versioning strict.

\subsection{Étape E16 --- Analyse des Contre-Expertises}

\textbf{Objectif~:} Évaluer les rapports contradictoires produits par la partie
adverse et préparer une réfutation méthodique.

\textbf{Actions~:} Analyse méthodologique critique~;
identification des erreurs factuelles~;
discussion des divergences interprétatives légitimes~;
rédaction d'un mémoire en réponse structuré~;
proposition d'expertises complémentaires si nécessaire.

\subsection{Étape E17 --- Clôture et Retour d'Expérience}

\textbf{Objectif~:} Archiver définitivement le dossier, capitaliser les
apprentissages et améliorer les processus futurs.

\textbf{Actions~:} Débriefing structuré d'équipe~;
rédaction d'un rapport de retour d'expérience (REX)~;
archivage définitif conforme aux durées légales~;
destruction sécurisée des données temporaires~;
intégration des enseignements dans la base de connaissances~;
mise à jour des procédures standard.

\begin{boiteRemarque}{E17 n'est pas optionnelle}
L'étape~E17 (\textit{Retour d'Expérience}) n'est pas optionnelle~: elle
constitue le mécanisme d'amélioration continue permettant l'évolution et
l'adaptation du processus aux nouveaux défis investigatifs. Une organisation
qui ne fait pas de REX après chaque investigation majeure est condamnée à
répéter les mêmes erreurs.
\end{boiteRemarque}

% ============================================================================
\section{Cas Pratique~: L'Affaire de la Facture Suspecte}
\label{sec:piu-rousseau}
% ============================================================================

\subsection{Contexte général}

Le 15~mars~2025, Madame Claire ROUSSEAU, dirigeante de la société INNOV-TECH
SARL (PME de 25~salariés), constate un débit de 47~850~\texteuro{} sur son relevé
bancaire au profit de la société CONSULTING-PRO, dont elle n'a aucune
connaissance. En vérifiant ses archives comptables, elle découvre une facture
n°~CP-2025-089 datée du 28~février~2025, pour des « prestations de conseil
stratégique »~: prestations qu'elle affirme n'avoir jamais commandées.

Cette investigation illustre comment le PIU s'applique différemment selon trois
contextes~:
\begin{enumerate}
    \item \textbf{Contexte particulier}~: Madame ROUSSEAU mène elle-même
          l'investigation~;
    \item \textbf{Contexte entreprise}~: INNOV-TECH mandate un audit interne~;
    \item \textbf{Contexte judiciaire}~: un expert est désigné par le tribunal.
\end{enumerate}

\subsection{Application~: Contexte particulier (niveau minimal/standard)}

\textbf{Phase~1~: Initialisation}

\textbf{E1~-- Réception et Évaluation~:} Enjeu financier~: 47~850 (significatif)~;
complexité~: moyenne~; faisabilité~: 10h/semaine disponibles~;
risque~: fraude interne ou externe. Décision~: investigation personnelle
préliminaire avant éventuel dépôt de plainte.

\textbf{E2~-- Décision d'Engagement~:} Niveau minimal~: décision seule.
Périmètre~: identifier l'origine de la facture et le responsable du paiement.
Délai~: 3~semaines. Budget~: temps personnel + consultation avocat (2h).
Pas de mandat formel, mais décision documentée dans un email à elle-même
(traçabilité).

\textbf{E3~-- Planification~:}
\begin{enumerate}
    \item Rassembler tous documents (facture, relevés, contrats,
          correspondances)~;
    \item Interroger ses employés (comptable, assistante administrative)~;
    \item Rechercher informations sur CONSULTING-PRO (KBIS, site web)~;
    \item Consulter avocat pour aspects juridiques.
\end{enumerate}
\textbf{QG1~:} Auto-évaluation VERT. Progression Phase~2.

\textbf{Phase~2~: Acquisition}

\textbf{E4~-- Collecte~:} Niveau standard~: facture CP-2025-089 (scan haute
résolution, horodaté)~; relevé bancaire~; extrait KBIS de CONSULTING-PRO~;
tous emails mentionnant « CONSULTING-PRO »~; registre courrier (pas de
mention)~; contrats 2024--2025 (aucun avec CONSULTING-PRO). Numérotation~:
EP001 à EP008.

\textbf{E5~-- Sécurisation~:} 3~copies~: disque dur externe chiffré (BitLocker),
cloud personnel sécurisé (ProtonDrive), copie papier en coffre. Hash SHA-256
calculé et conservé séparément pour chaque fichier numérique. Règle~: jamais
travailler sur originaux papier.

\textbf{E6~-- Parties Prenantes~:} Sophie LEMOINE (comptable, a enregistré
la facture)~; Julie MARTIN (assistante administrative)~; Marc DUBOIS (gérant
CONSULTING-PRO selon KBIS)~; Me Paul BERNARD (avocat conseil).

\textbf{QG2~:} 8~éléments numérotés~; 3~copies~; hash calculés~; journal
d'accès manuscrit. VERT. Progression Phase~3.

\textbf{Phase~3~: Investigation}

\textbf{E7~-- Interrogatoires~:} Sophie LEMOINE (12~mars, 14h)~: déclare avoir
reçu la facture par email le 1er~mars~; semblait régulière, tamponnée
« bon à payer »~; ne connaît pas CONSULTING-PRO~; fournit l'email de réception
(\texttt{contact@consulting-pro.fr}). Julie MARTIN (13~mars)~: aucune trace.
Comptes rendus signés.

\textbf{E8~-- Hypothèses~:}
\begin{itemize}
    \item \textbf{H1}~: Erreur administrative --- facture légitime oubliée~;
    \item \textbf{H2}~: Fraude externe --- CONSULTING-PRO frauduleuse~;
    \item \textbf{H3}~: Complicité interne --- Sophie complice~;
    \item \textbf{H4}~: Usurpation d'identité entreprise~;
    \item \textbf{H5}~: Erreur de destinataire (facture destinée à homonyme).
\end{itemize}

\textbf{E9~-- Confrontation~:} H1 réfutée (gestion rigoureuse, aucun contrat)~;
H2 cohérente (KBIS~: création 15~janvier~2025, très récent, suspect)~;
H3 peu cohérente (Sophie sincère, 5~ans ancienneté, pas d'indice financier)~;
H4 possible mais sans preuve~; H5 réfutée (facture mentionne INNOV-TECH
avec SIRET correct). Score~: H2 dominant.

\textbf{E9bis~-- Falsification~:} Pour H2~: si prestation réellement effectuée,
il devrait exister~: contrat signé, correspondances préalables, livrables.
Recherche approfondie~: rien trouvé. H2 renforcée.

\textbf{E10~-- Validation~:} Consultation de Me BERNARD (avocat)~: confirme
indices sérieux de fraude~; recommande dépôt de plainte.

\textbf{QG3~:} 5~hypothèses générées~; matrice complétée~; falsification
effectuée~; validation externe (avocat). VERT.

\textbf{Phase~4~: Formalisation}

\textbf{E11--E12~-- Rapport~:} Dossier simple (15~pages)~: synthèse
(facture suspecte, aucune prestation commandée, indices de fraude externe)~;
chronologie~; copies EP001--EP008~; comptes rendus d'entretiens~;
conclusion~: fraude externe présumée (H2). Transmission~: avocat (chiffré)
+ commissariat (copie papier).

\textbf{Phase~5~:} Maintien du dossier pour enquête policière éventuelle.
Documentation de toutes les suites.

\subsection{Application~: Contexte entreprise (niveau avancé)}

Le 16~mars, Madame ROUSSEAU alerte son conseil d'administration. Décision
collégiale~: audit interne approfondi. Mandat formel signé par le Président
du CA. Équipe~: Jean DUPUIS (Dir.~Financier, lead)~+ Marie LECLERC (DRH)
+ Me Sophie MOREAU (Avocate). Budget~: 15~000. Délai~: 6~semaines.

\textbf{Phase~2~-- Niveau avancé~:} Image forensique du serveur de messagerie
(période janvier--mars~2025)~; extraction logs accès ERP comptable~;
OSINT sur CONSULTING-PRO (révèle antécédents judiciaires de Marc DUBOIS~:
escroquerie 2019)~; présence huissier pour constat de l'email suspect~;
hash SHA-512 sur toutes images forensiques~; blockchain privée pour
horodatage immuable.

\textbf{Phase~3~-- Niveau avancé~:} Expertise IT sur l'email révèle~: spoofing
(falsification de l'adresse expéditeur)~; headers techniques montrent origine
depuis serveur russe~; PDF de la facture créé avec logiciel de retouche.
Processus ACH formalisé~: 7~hypothèses, matrice quantitative bayésienne,
tests de falsification systématiques. Conclusion~: H2 confirmée avec
probabilité~$> 95\%$. Red Team interne~: consensus 100\%.

\textbf{Phase~4~:} Rapport 85~pages en 3~volumes~; présentation orale au
CA~; archivage légal 10~ans.

\subsection{Application~: Contexte judiciaire (niveau avancé systématique)}

Le 1er~avril~2025, le juge d'instruction désigne Madame Isabelle FONTAINE,
expert judiciaire, pour répondre à deux questions~: (1)~la facture correspond-elle
à une prestation réelle~? (2)~quelle est son origine~?

\textbf{Phase~1~:} Vérification d'absence de conflit d'intérêt~; serment prêté.
Mandat conforme~: ordonnance~; méthodologie~: \norme{ISO~27037} pour forensique,
principes CNCEJ pour expertise judiciaire~; délai~: 4~mois.

\textbf{Phase~2~:} Collecte contradictoire en présence des parties, avec
greffier. Saisie contradictoire du serveur messagerie INNOV-TECH (15~avril)~;
convocation de Marc DUBOIS (22~avril)~: refuse de produire tout document~;
scellés judiciaires~; hash SHA-512 certifié par l'expert.

\textbf{Phase~3~:} ACH formalisée avec documentation exhaustive~;
conclusion~: H2 (fraude externe par phishing/spoofing) probabilité~$> 98\%$~;
peer review par confrère expert.

\textbf{Phase~4~:} Rapport d'expertise judiciaire (65~pages, normes CNCEJ)~:
réponses aux questions du juge~:}
\begin{itemize}
    \item \textit{Q1 -- La facture correspond-elle à une prestation réelle~?}
          \textbf{R~: Non. Aucun élément probatoire ne démontre qu'une
          prestation ait été effectuée.}
    \item \textit{Q2 -- Origine de la facture~?}
          \textbf{R~: Émission frauduleuse par CONSULTING-PRO via
          techniques de phishing/spoofing.}
\end{itemize}
Dépôt au greffe le 25~juillet~; archivage 30~ans.

\textbf{Phase~5~:} Contre-expertise privée de CONSULTING-PRO (prétend prestation
orale)~: réfutée par l'expert (une prestation de 47k ne peut être exclusivement
orale). Condamnation pénale de Marc DUBOIS (escroquerie).

\subsection{Synthèse comparative des trois contextes}

\begin{tableaucours}{p{2.8cm}p{3.2cm}p{3.2cm}p{3.2cm}}{%
    \tableauHeader{Dimension} &
    \tableauHeader{Particulier} &
    \tableauHeader{Entreprise} &
    \tableauHeader{Judiciaire}
}
Niveau rigueur & Minimal/Standard & Standard/Avancé & Avancé systématique \\
Durée & 3~semaines & 6~semaines & 4~mois \\
Ressources & 1~personne + avocat (2h) & Équipe~3 + outils forensiques &
    Expert + consultant IT \\
Coût estimé & 500~(avocat) & 15~000 & 8~000~(honoraires) \\
Validation & Avocat & Red Team interne & Peer review + juge \\
Livrable & Dossier 15p + plainte & Rapport 85p (3~vol.) &
    Rapport expertise 65p (CNCEJ) \\
Force probante & Moyenne & Élevée (interne) & Maximale (judiciaire) \\
\end{tableaucours}
\vspace{-0.8em}
\begin{center}{\small\itshape Tableau~6.1 --- Comparaison des trois applications
contextuelles du PIU}\end{center}

% ============================================================================
\section{Recommandations d'Usage et Domaines d'Application}
\label{sec:piu-recommandations}
% ============================================================================

\subsection{Principe fondamental}

La qualité d'une investigation se mesure moins à la rapidité avec laquelle on
parvient à des conclusions qu'à la robustesse avec laquelle ces conclusions
résistent à la contestation contradictoire. Le PIU, par sa structure formelle
intégrant falsification active, validation multi-méthodes et Quality Gates,
vise cet objectif de défendabilité maximale.

\subsection{Sélection du niveau de rigueur}

Le praticien doit calibrer le niveau de rigueur selon la matrice de criticité
(Tableau~\ref{tab:piu-criticite}). Une investigation à enjeu limité ne justifie
pas systématiquement un niveau avancé (coût/bénéfice)~; une investigation à
enjeu critique (pénal, financier majeur) ne peut se satisfaire d'un niveau
minimal sans risquer l'invalidation ultérieure.

\subsection{Domaines d'application}

Bien que ce cours illustre le PIU dans un contexte de forensique numérique,
la méthodologie s'applique à tout domaine investigatif~:
\begin{itemize}
    \item Investigations RH (harcèlement, discrimination)~;
    \item Forensique numérique (cybercriminalité, fuites de données)~;
    \item Investigations techniques (accidents, défaillances)~;
    \item Investigations scientifiques (fraude académique)~;
    \item Investigations journalistiques~;
    \item Audits de conformité réglementaire.
\end{itemize}

\subsection{Documentation comme investissement}

L'effort documentaire demandé à chaque étape (inventaires, matrices, journaux)
peut sembler lourd mais constitue un \textbf{investissement essentiel}.
En cas de contestation (Phase~5), la qualité de la documentation initiale
détermine directement la capacité à défendre les conclusions, parfois plusieurs
années après les faits.

\begin{boxCRO}
\textbf{Analyse CRO --- Le PIU dans son ensemble}

Le PIU est conçu pour maximiser $\mathcal{O}$ (opposabilité) tout en
maintenant $\mathcal{R}$ (fiabilité) et en gérant $\mathcal{C}$ (confidentialité).

Les Quality Gates sont les mécanismes d'assurance qualité qui garantissent
que $\mathcal{R}$ et $\mathcal{O}$ ne sont pas compromises par la pression
temporelle ou le biais de confirmation. La falsification active (E9bis) est
le dispositif épistémologique qui protège $\mathcal{R}$ contre les biais.

La confidentialité ($\mathcal{C}$) est gérée par les niveaux de rigueur~:
le niveau avancé (contexte judiciaire) exige une protection maximale des
données sensibles~; le niveau minimal (contexte particulier) accepte
davantage de compromis proportionnels à l'enjeu.

\cro{$\sim$ gérée par niveau}{$\uparrow\uparrow$ maximisée}{$\uparrow\uparrow$ objectif principal}
\end{boxCRO}

\separateur

\begin{boxresume}
\textbf{Ce qu'il faut retenir de ce chapitre~:}
\begin{itemize}
    \item Le \textbf{PIU} est un processus modulaire en \textbf{5~phases,
          17~étapes et 3~Quality Gates}~: Initialisation (E1--E3) $\to$
          Acquisition (E4--E6) $\to$ Investigation (E7--E10) $\to$
          Formalisation (E11--E13) $\to$ Suivi (E14--E17).

    \item La \textbf{matrice de criticité} calibre le niveau de rigueur (minimal
          / standard / avancé) selon l'enjeu financier, pénal, technique,
          le nombre de parties et la durée.

    \item Les \textbf{Quality Gates} (QG1~: pré-opérationnel~;
          QG2~: intégrité probatoire~; QG3~: robustesse analytique) sont
          des points de contrôle obligatoires à trois niveaux~:
          VERT (progression)~/ ORANGE (conditionnel)~/ ROUGE (suspension).

    \item La \textbf{boucle itérative} entre Phase~2 et Phase~3 est normale~:
          2 à 3 passages est satisfaisant~; un passage direct suggère un
          manque d'approfondissement.

    \item La \textbf{falsification active} (E9bis) -- principe poppérien --
          est la protection contre le biais de confirmation. Chercher à
          réfuter ses hypothèses, pas à les confirmer.

    \item L'\textbf{affaire ROUSSEAU} illustre que le \emph{même} fait
          (fraude externe H2 à 47~850) est atteint par trois niveaux de
          rigueur, mais avec une force probante radicalement différente~:
          suffisante pour une plainte, élevée pour une décision interne,
          maximale pour une condamnation pénale.
\end{itemize}
\end{boxresume}

\bigskip

\begin{boxexo}[title={\ding{45}\quad Exercice~6.1 --- Matrice de criticité \niveaubase}]
Pour chacune des investigations suivantes, déterminez le niveau de rigueur
approprié (minimal / standard / avancé) en justifiant votre réponse par
la matrice de criticité~:
\begin{enumerate}[(a)]
    \item Soupçon de fraude interne de 450~000~FCFA dans une administration
          publique camerounaise~; suspect unique~; durée estimée~: 2~mois.
    \item Recherche de documents compromettants sur le PC d'un employé
          licencié pour harcèlement moral~; aucune plainte déposée~;
          enjeu~: décision de licenciement interne.
    \item Investigation mandatée par le procureur sur une attaque ransomware
          ayant paralysé un hôpital pendant 72h~; données patients exposées~;
          plusieurs suspects identifiés dans 3~pays.
\end{enumerate}
\end{boxexo}

\begin{boxexo}[title={\ding{45}\quad Exercice~6.2 --- Application du PIU \niveaumoyen}]
Vous êtes mandaté par la direction de TRANSBOIS~SA (affaire MVOM,
Chapitre~\ref{chap:gr-affaires}) pour conduire l'investigation interne
sur l'attaque ransomware. Appliquez les étapes E1 à E6~:
\begin{enumerate}[(a)]
    \item Rédigez la \textit{Fiche d'Évaluation Initiale} (E1)~: enjeux,
          risques, faisabilité.
    \item Définissez le périmètre d'investigation (E2) et justifiez le niveau
          de rigueur choisi.
    \item Listez les 8~premières sources à collecter (E4) en les classant
          par ordre de volatilité~; précisez la méthode de collecte pour
          chacune.
    \item Quels critères QG2 risquent d'être difficiles à satisfaire dans
          ce contexte~? Proposez comment les atteindre.
\end{enumerate}
\end{boxexo}

\begin{boxexo}[title={\ding{45}\quad Exercice~6.3 --- Falsification active \niveauexpert}]
Dans l'investigation ROUSSEAU (contexte entreprise), l'hypothèse H2
(fraude externe par phishing/spoofing) a été retenue avec une probabilité
$> 95\%$.

\begin{enumerate}[(a)]
    \item Identifiez 4~observations critiques qui, si trouvées, auraient
          définitivement réfuté H2 (conditions de falsification).
    \item Comment calculer l'indice de robustesse de H2~? Quel score
          attribuez-vous à cette hypothèse~?
    \item Formulez l'analyse CRO de la décision d'inclure dans le rapport
          l'analyse forensique des en-têtes email (révèle l'origine depuis
          un serveur russe, potentiellement sensible géopolitiquement).
    \item Dans quel cas aurait-il fallu activer le niveau ROUGE au QG3~?
          Quelles étapes de retour auraient été nécessaires~?
\end{enumerate}
\end{boxexo}

\bigskip

\begin{boxinfo}
\textbf{Bibliographie recommandée pour ce chapitre~:}
\begin{itemize}
    \item \norme{ISO/IEC~27037:2012}~: Guidelines for identification,
          collection, acquisition and preservation of digital evidence.
    \item \norme{NIST SP~800-86}~: Guide to Integrating Forensic Techniques
          into Incident Response, 2006.
    \item Heuer, R.~J.,
          \textit{Psychology of Intelligence Analysis},
          CIA, 1999 (méthode ACH). \cite{heuer1999}
    \item Popper, K., \textit{La logique de la découverte scientifique},
          Payot, 1973. \cite{popper1973}
    \item ACPO, \textit{Good Practice Guide for Digital Evidence}, 2012.
\end{itemize}
\textbf{Transition.} Le Chapitre~\ref{chap:meth-std} compare le PIU aux
standards internationaux (SANS FOR508, NIST, DFRWS, ACPO) pour en
montrer la cohérence et les spécificités. Les Chapitres~\ref{chap:custody}
et~\ref{chap:acq} approfondissent respectivement la chaîne de custody
et les outils d'acquisition --- deux dimensions essentielles de la Phase~2.
\end{boxinfo}